\chapter{Advanced Model Features}
\label{chap:advanced-features}

This chapter collects model features that extend the basic mesh--material--load workflow. Concentrated inertia has two supported input representations: native \cmd{POINTMASS} assigns generalized point properties directly to an NSET, while \cmd{ELEMENT, TYPE=MASS} together with \cmd{MASS} provides the shared FEMaster/Abaqus one-node element/property representation. Both paths use the same point-element and point-mass-section mechanics internally. The topology-static commands use element Fields to scale stiffness and rotate material directions without changing mesh connectivity.

\keywordpage{POINTMASS}

\cmd{POINTMASS} is the native NSET-based concentrated-property command. It is evaluated on the compiled model and creates one auxiliary point element for every node in the target NSET. All point elements created by one command share the supplied mass, rotary inertia, and optional ground-spring properties.

The first value is one translational point mass \(m\). It contributes equally to the three translational inertia directions. The next three values are rotary inertias
\[
(I_x,I_y,I_z)
\]
associated with the nodal rotational degrees of freedom. Translational and rotational spring constants may then be supplied. Missing values default to zero, so a pure lumped mass requires only the first scalar.

For one target node, the diagonal generalized contributions are
\[
M_p=\operatorname{diag}(m,m,m,I_x,I_y,I_z),
\]
and
\[
K_p=\operatorname{diag}(k_x,k_y,k_z,k_{rx},k_{ry},k_{rz}).
\]
All coefficients must be non-negative.

The point element activates only generalized components that carry a nonzero physical coefficient. A nonzero translational mass activates \(U_x,U_y,U_z\); a nonzero rotary inertia activates its corresponding rotational component; translational and rotational springs likewise activate their associated directions. A point element with all coefficients zero is mechanically inactive.

If the target set contains several nodes, the complete property is applied independently at every target node. The specified mass is therefore mass \emph{per node}, not the total mass of the NSET. A total payload distributed over four target nodes requires four appropriately chosen nodal masses whose sum equals the intended total.

The auxiliary point elements created by \cmd{POINTMASS} are deliberately not inserted into the public compiled element array or any ELSET. This preserves ordinary element numbering and element-domain Field offsets. They nevertheless participate in active-DOF enumeration, global stiffness assembly, global mass assembly, eigenfrequency and transient analyses, and the dedicated point-mass-aware inertia paths.

This distinction matters for loads. Native \cmd{VLOAD} and supported Abaqus \cmd{DLOAD, GRAV} operate on ordinary element regions and therefore do not select auxiliary \cmd{POINTMASS} objects. Native \cmd{INERTIALOAD} includes them only when \key{CONSIDER_POINT_MASSES=1}. Likewise, native \cmd{INERTIARELIEF} includes them in total mass, center of gravity, and rotary inertia only when its \key{CONSIDER_POINT_MASSES} option is enabled.

Because the auxiliary point elements activate their own DOFs, a model may consist entirely of nodes plus \cmd{POINTMASS} definitions without requiring an unrelated structural element merely to activate translational or rotational components. Node-only models can also be written to FRD; the writer omits the element-connectivity block when no regular output element topology exists.

Point masses represent equipment, payloads, actuators, engines, sensors, or other concentrated inertia without explicitly meshing the attached body. Translational stiffness has force/length units, rotational stiffness has moment/radian units, and rotary inertia has mass-length\(^2\) units in the chosen consistent system.

\begin{keywordparameters}
\key{NSET} & required & Target compiled node set. \\
\end{keywordparameters}

\begin{keyworddata}
\val{mass} & Concentrated translational mass; default 0. \\
\val{Ix, Iy, Iz} & Rotary inertia components; default 0. \\
\val{kx, ky, kz} & Translational ground-spring constants; default 0. \\
\val{krx, kry, krz} & Rotational ground-spring constants; default 0. \\
\end{keyworddata}

\exampleheading
\begin{dialectexamples}
\begin{femasterdeck}
*POINTMASS, NSET=PAYLOAD
0.025,
1.2e-3, 1.2e-3, 2.0e-3,
1000., 1000., 0.,
0., 0., 0.
\end{femasterdeck}
\begin{noabaqus}
Abaqus-style input uses \cmd{ELEMENT, TYPE=MASS} together with \cmd{MASS} for the supported isotropic translational mass. Rotary inertia and ground-spring terms from native \cmd{POINTMASS} are not provided by that mapping.
\end{noabaqus}
\end{dialectexamples}

\keywordpage{MASS}

\cmd{MASS} assigns an isotropic concentrated translational mass to regular one-node elements created with \cmd{ELEMENT, TYPE=MASS}. \key{ELSET} identifies the target point-element region. The command is accepted by both the native and supported Abaqus readers.

The topology and physical property are deliberately separate. A \val{MASS} element stores one element ID, one connected node, and normal ELSET membership. Without a property it is inert: it has zero mass, zero stiffness, no stress/strain state, and activates no DOFs. \cmd{MASS} attaches a point-mass section whose supported imported form contains one scalar \(m\). The resulting element mass matrix is
\[
M_e=\operatorname{diag}(m,m,m,0,0,0),
\]
so a positive mass activates the three translational nodal directions.

The supplied scalar is a property of every target element. If an ELSET contains four distinct \val{MASS} elements and the input value is \val{0.01}, the total represented translational mass is \val{0.04} in the chosen consistent unit system. Several MASS elements may reference the same node; their assembled mass contributions remain independent and additive. The complete target ELSET is validated before assignment and mixed sets containing non-MASS elements are rejected.

Regular \val{MASS} elements participate in the ordinary structural-element paths. They contribute to the global mass matrix, eigenfrequency and transient response, and inertia-relief mass properties. Their density-scaled field integration treats the section mass as the complete concentrated measure, so a native \cmd{VLOAD} or supported Abaqus \cmd{DLOAD, GRAV} produces the equivalent concentrated force \(m\boldsymbol a\) whenever the selected element region contains the MASS element. Native \cmd{INERTIALOAD} likewise includes regular MASS elements when they belong to its target ELSET; the \key{CONSIDER_POINT_MASSES} switch applies only to auxiliary native \cmd{POINTMASS} objects.

Root- and Part-level \cmd{MASS} assignments are stored with their semantic topology before compilation. Part-local assignments therefore follow every Instance of that Part automatically. Assembly-level \cmd{MASS} records are resolved after compilation against assembly ELSETs and can target selected placed point elements directly.

\key{TYPE} defaults to \val{ISOTROPIC} and currently accepts only that value. Abaqus anisotropic mass, orientation-dependent mass, rotary inertia, and damping parameters are not mapped by this command. Native \cmd{POINTMASS} remains the more general input when rotary inertia or ground-spring stiffness is required.

\begin{keywordparameters}
\key{ELSET} & required & Nonempty element set containing only \val{TYPE=MASS} point elements. \\
\key{TYPE} & optional & \val{ISOTROPIC}; currently the only supported form and the default. \\
\end{keywordparameters}

\begin{keyworddata}
\val{mass} & Non-negative isotropic translational mass assigned to every target point element. \\
\end{keyworddata}

\exampleheading
\begin{dialectexamples}
\begin{femasterdeck}
*ELEMENT, TYPE=MASS, ELSET=PAYLOAD
100, 2
*MASS, ELSET=PAYLOAD
0.01
\end{femasterdeck}
\begin{abaqusdeck}
*ELEMENT, TYPE=MASS, ELSET=PAYLOAD
100, 2
*MASS, ELSET=PAYLOAD
0.01
\end{abaqusdeck}
\end{dialectexamples}

\section{Topology-weighted linear statics}

\val{LINEARSTATICTOPO} extends linear static analysis with element-wise density and optional material-orientation Fields. It is intended for topology/sizing workflows in which design variables are supplied externally or updated between solver runs.

The density Field scales each element's constitutive stiffness by
\[
s_e=\rho_e^p,
\]
where \(\rho_e\) is the value selected by \cmd{TOPODENSITY} and \(p\) is set by \cmd{TOPOEXPONENT}. Thus an element with \(\rho=1\) retains its full stiffness, while intermediate density is penalized according to the exponent.

If an orientation Field is supplied, its three components are interpreted as element-wise Euler rotation angles. The resulting additional rotation is applied to the material coordinate system for the topology-aware solid response. The Field therefore changes material orientation, not element geometry.

The analysis also produces topology-related output such as density, compliance-related quantities, density sensitivity, volume, and orientation sensitivity where applicable.

\keywordpage{TOPODENSITY}

\cmd{TOPODENSITY} selects the scalar element Field \(\rho_e\) used by an active \val{LINEARSTATICTOPO} loadcase. \key{FIELD} is required; \key{NAME} is accepted as an alternative spelling. The referenced Field must use \val{TYPE=ELEMENT} and exactly one component.

The stiffness multiplier is
\[
\rho_e^p.
\]
The command does not automatically force values into the interval \([0,1]\). An optimization based on physical density fractions therefore depends on an external workflow that supplies values consistent with that interpretation. Negative density values or inappropriate exponents can produce nonphysical stiffness scaling.

For a SIMP-style topology problem, \(0<\rho_\mathrm{min}\leq\rho_e\leq1\) is commonly used to avoid exactly singular void elements, but the choice of minimum density belongs to the optimization strategy surrounding the FEM solve.

Compliance sensitivity contains factors proportional to \(1/\rho_e\). Exactly zero density therefore introduces singular sensitivity behaviour in gradient-based topology workflows even though the associated zero-stiffness element has a direct conceptual interpretation.

\begin{keywordparameters}
\key{FIELD} / \key{NAME} & required & Scalar \val{ELEMENT} Field containing element densities. \\
\end{keywordparameters}

\exampleheading
\begin{dialectexamples}
\begin{femasterdeck}
*FIELD, NAME=RHO, TYPE=ELEMENT,
 COLS=1, FILL=ZERO
0, 0.8
1, 1.0

*LOADCASE, TYPE=LINEARSTATICTOPO
*TOPODENSITY, FIELD=RHO
*TOPOEXPONENT
3.0
*END
\end{femasterdeck}
\begin{noabaqus}
Topology-density Fields and \val{LINEARSTATICTOPO} are native FEMaster capabilities.
\end{noabaqus}
\end{dialectexamples}

\keywordpage{TOPOORIENT}

\cmd{TOPOORIENT} selects a three-component \val{ELEMENT} Field containing additional material-orientation angles for \val{LINEARSTATICTOPO}. \key{FIELD} is required and \key{NAME} is accepted as an alias.

The three Field components are Euler rotation angles
\[
(\theta_x,\theta_y,\theta_z)
\]
in \textbf{radians}. They define the additional material rotation
\[
R(\theta_x,\theta_y,\theta_z)
=R_x(\theta_x)\,R_y(\theta_y)\,R_z(\theta_z).
\]
This rotation acts in addition to the Section's base material orientation. It is therefore intended for optimization or parameter studies in which the preferred anisotropic material direction varies from element to element.

An orientation Field changes only directional constitutive response. It does not rotate the mesh and has no effect on a fully isotropic material law because isotropic stiffness is invariant under rotation.

The rotation order is significant because finite rotations about different axes do not commute. An external optimization or preprocessing tool that generates \cmd{TOPOORIENT} values must therefore use the same \(R_xR_yR_z\) convention and radians. For example, a second-component value of \val{0.2} means a rotation of \(0.2\,\mathrm{rad}\), not \(0.2^\circ\).

\begin{keywordparameters}
\key{FIELD} / \key{NAME} & required & Three-component \val{ELEMENT} Field of Euler angles \((\theta_x,\theta_y,\theta_z)\) in radians. \\
\end{keywordparameters}

\exampleheading
\begin{dialectexamples}
\begin{femasterdeck}
*FIELD, NAME=ORIENT,
 TYPE=ELEMENT, COLS=3, FILL=ZERO
0, 0.0, 0.0, 0.0
1, 0.0, 0.2, 0.0

*LOADCASE, TYPE=LINEARSTATICTOPO
*TOPODENSITY, FIELD=RHO
*TOPOORIENT, FIELD=ORIENT
*END
\end{femasterdeck}
\begin{noabaqus}
No topology-orientation Field mapping is currently supported by the Abaqus reader.
\end{noabaqus}
\end{dialectexamples}

\keywordpage{TOPOEXPONENT}

\cmd{TOPOEXPONENT} sets the scalar penalization exponent \(p\) used by \val{LINEARSTATICTOPO}:
\[
K_e(\rho_e)=\rho_e^p K_{e,0}
\]
for the element stiffness contribution represented by the topology scaling.

With \(p=1\), stiffness varies linearly with density. Values \(p>1\) penalize intermediate densities increasingly strongly and are commonly used in SIMP-style topology optimization to drive the design toward a more nearly discrete solid/void distribution.

The exponent also affects compliance sensitivity. If an element compliance contribution before the density derivative is denoted by \(c_e\), the density sensitivity scales conceptually with
\[
\frac{\partial C}{\partial\rho_e}
\propto -p\,\rho_e^{p-1}c_e.
\]
The exact reported Field follows FEMaster's topology result definitions, while the sign and penalization dependence follow this standard stiffness-scaling idea.

\begin{keyworddata}
\val{exponent} & Scalar density penalization exponent \(p\). \\
\end{keyworddata}

\exampleheading
\begin{dialectexamples}
\begin{femasterdeck}
*TOPOEXPONENT
3.0
\end{femasterdeck}
\begin{noabaqus}
No corresponding topology penalization keyword is currently supported in Abaqus input.
\end{noabaqus}
\end{dialectexamples}
